\Chapter{77}

\Verse{1}
{
    \zA{话说王夫人见中秋已过，凤姐病也比先减了，虽未大愈，然亦可以出入行走得 了，仍命大夫每日诊脉服药。}
}
{
    \eA{[English source not present in the checked-in Bencraft/Joly files.]}
}
{
}

\Verse{2}
{
    \zA{又开了丸药方来，配“调经养荣丸”。}
}
{
    \eA{[English source not present in the checked-in Bencraft/Joly files.]}
}
{
}

\Verse{3}
{
    \zA{因用上等人参 二两，王夫人取时，翻寻了半日，只向小匣内寻了几枝簪挺粗细的。}
}
{
    \eA{[English source not present in the checked-in Bencraft/Joly files.]}
}
{
}

\Verse{4}
{
    \zA{王夫人看了嫌 不好，命再找去，又找了一大包须沫出来。}
}
{
    \eA{[English source not present in the checked-in Bencraft/Joly files.]}
}
{
}

\Verse{5}
{
    \zA{王夫人焦躁道：“用不着偏有，但用着 了，再找不着!成日家我叫你们查一查，都归拢一处，你们白不听，就随手混撂。”}
}
{
    \eA{[English source not present in the checked-in Bencraft/Joly files.]}
}
{
}

\Verse{6}
{
    \zA{彩云道：“想是没了，就只有这个。}
}
{
    \eA{[English source not present in the checked-in Bencraft/Joly files.]}
}
{
}

\Verse{7}
{
    \zA{上次那边的太太来寻了去了。”}
}
{
    \eA{[English source not present in the checked-in Bencraft/Joly files.]}
}
{
}

\Verse{8}
{
    \zA{王夫人道：“没 有的话。}
}
{
    \eA{[English source not present in the checked-in Bencraft/Joly files.]}
}
{
}

\Verse{9}
{
    \zA{你再细找找。”}
}
{
    \eA{[English source not present in the checked-in Bencraft/Joly files.]}
}
{
}

\Verse{10}
{
    \zA{彩云只得又去找寻，拿了几包药材来，说：“我们不认的 这个，请太太自看。}
}
{
    \eA{[English source not present in the checked-in Bencraft/Joly files.]}
}
{
}

\Verse{11}
{
    \zA{除了这个没有了。”}
}
{
    \eA{[English source not present in the checked-in Bencraft/Joly files.]}
}
{
}

\Verse{12}
{
    \zA{王夫人打开看时，也都忘了，不知都是什 么，并没有一支人参。}
}
{
    \eA{[English source not present in the checked-in Bencraft/Joly files.]}
}
{
}

\Verse{13}
{
    \zA{因一面遣人去问凤姐有无。}
}
{
    \eA{[English source not present in the checked-in Bencraft/Joly files.]}
}
{
}

\Verse{14}
{
    \zA{凤姐来说：“也只有些参膏。}
}
{
    \eA{[English source not present in the checked-in Bencraft/Joly files.]}
}
{
}

\Verse{15}
{
    \zA{芦 须虽有几根，也不是上好的，每日还要煎药里用呢。”}
}
{
    \eA{[English source not present in the checked-in Bencraft/Joly files.]}
}
{
}

\Verse{16}
{
    \zA{王夫人听了，只得向邢夫人 那里问去。}
}
{
    \eA{[English source not present in the checked-in Bencraft/Joly files.]}
}
{
}

\Verse{17}
{
    \zA{说：“因上次没了，才往这里来寻，早已用完了。”}
}
{
    \eA{[English source not present in the checked-in Bencraft/Joly files.]}
}
{
}

\Verse{18}
{
    \zA{王夫人没法，只得 亲身过来请问贾母。}
}
{
    \eA{[English source not present in the checked-in Bencraft/Joly files.]}
}
{
}

\Verse{19}
{
    \zA{贾母忙命鸳鸯取出当日馀的来，竟还有一大包，皆有手指头粗 细不等，遂秤了二两给王夫人。}
}
{
    \eA{[English source not present in the checked-in Bencraft/Joly files.]}
}
{
}

\Verse{20}
{
    \zA{王夫人出来，交给周瑞家的拿去，令小厮送与医生 家去。}
}
{
    \eA{[English source not present in the checked-in Bencraft/Joly files.]}
}
{
}

\Verse{21}
{
    \zA{又命将那几包不能辨的药也带了去，命医生认了，各包号上。}
}
{
    \eA{[English source not present in the checked-in Bencraft/Joly files.]}
}
{
}

\Verse{22}
{
    \zA{一时周瑞家的 又拿进来，说：“这几样都各包号上名字了。}
}
{
    \eA{[English source not present in the checked-in Bencraft/Joly files.]}
}
{
}

\Verse{23}
{
    \zA{但那一包人参固然是上好的，只是年 代太陈。}
}
{
    \eA{[English source not present in the checked-in Bencraft/Joly files.]}
}
{
}

\Verse{24}
{
    \zA{这东西比别的却不同，凭是怎么好的，只过一百年后，就自己成了灰了。}
}
{
    \eA{[English source not present in the checked-in Bencraft/Joly files.]}
}
{
}

\Verse{25}
{
    \zA{如今这个虽未成灰，然已成了糟朽烂木，也没有力量的了。}
}
{
    \eA{[English source not present in the checked-in Bencraft/Joly files.]}
}
{
}

\Verse{26}
{
    \zA{请太太收了这个，倒不 拘粗细，多少再换些新的才好。”}
}
{
    \eA{[English source not present in the checked-in Bencraft/Joly files.]}
}
{
}

\Verse{27}
{
    \zA{王夫人听了，低头不语，半日才说：“这可没法了，只好去买二两来罢。”}
}
{
    \eA{[English source not present in the checked-in Bencraft/Joly files.]}
}
{
}

\Verse{28}
{
    \zA{也 无心看那些，只命：“都收了罢。”}
}
{
    \eA{[English source not present in the checked-in Bencraft/Joly files.]}
}
{
}

\Verse{29}
{
    \zA{因问周瑞家的：“你就去说给外头人们，拣好 的换二两来。}
}
{
    \eA{[English source not present in the checked-in Bencraft/Joly files.]}
}
{
}

\Verse{30}
{
    \zA{倘或一时老太太问你们，只说用的是老太太的，不必多说。”}
}
{
    \eA{[English source not present in the checked-in Bencraft/Joly files.]}
}
{
}

\Verse{31}
{
    \zA{周瑞家 的方才要去时，宝钗因在坐，乃笑道：“姨娘且住，如今外头人参都没有好的。}
}
{
    \eA{[English source not present in the checked-in Bencraft/Joly files.]}
}
{
}

\Verse{32}
{
    \zA{虽 有全枝，他们也必截做两三段，镶嵌上芦泡须枝，搀匀了好卖，看不得粗细。}
}
{
    \eA{[English source not present in the checked-in Bencraft/Joly files.]}
}
{
}

\Verse{33}
{
    \zA{我们 铺子里常和行里交易，如今我去和妈妈说了，哥哥去托个伙计过去和参行里要他二 两原枝来，不妨咱们多使几两银子，到底得了好的。”}
}
{
    \eA{[English source not present in the checked-in Bencraft/Joly files.]}
}
{
}

\Verse{34}
{
    \zA{王夫人笑道：“倒是你明白。}
}
{
    \eA{[English source not present in the checked-in Bencraft/Joly files.]}
}
{
}

\Verse{35}
{
    \zA{但只还得你亲自走一趟，才能明白。”}
}
{
    \eA{[English source not present in the checked-in Bencraft/Joly files.]}
}
{
}

\Verse{36}
{
    \zA{于是宝钗去了，半日回来说：“已遣人去， 赶晚就有回信。}
}
{
    \eA{[English source not present in the checked-in Bencraft/Joly files.]}
}
{
}

\Verse{37}
{
    \zA{明日一早去配也不迟。”}
}
{
    \eA{[English source not present in the checked-in Bencraft/Joly files.]}
}
{
}

\Verse{38}
{
    \zA{王夫人自是喜悦，因说道：“‘卖油的娘 子水梳头’。}
}
{
    \eA{[English source not present in the checked-in Bencraft/Joly files.]}
}
{
}

\Verse{39}
{
    \zA{自来家里有的给人多少，这会子轮到自己用，反倒各处寻去。”}
}
{
    \eA{[English source not present in the checked-in Bencraft/Joly files.]}
}
{
}

\Verse{40}
{
    \zA{说毕 长叹。}
}
{
    \eA{[English source not present in the checked-in Bencraft/Joly files.]}
}
{
}

\Verse{41}
{
    \zA{宝钗笑道：“这东西虽然值钱，总不过是药，原该济众散人才是。}
}
{
    \eA{[English source not present in the checked-in Bencraft/Joly files.]}
}
{
}

\Verse{42}
{
    \zA{咱们比不 得那没见世面的人家，得了这个，就珍藏密敛的。”}
}
{
    \eA{[English source not present in the checked-in Bencraft/Joly files.]}
}
{
}

\Verse{43}
{
    \zA{王夫人点头道：“你这话也是。”}
}
{
    \eA{[English source not present in the checked-in Bencraft/Joly files.]}
}
{
}

\Verse{44}
{
    \zA{一时宝钗去后，因见无别人在室，遂唤周瑞家的，问：“前日园中搜检的事情， 可得下落？”}
}
{
    \eA{[English source not present in the checked-in Bencraft/Joly files.]}
}
{
}

\Verse{45}
{
    \zA{周瑞家的是已和凤姐商议停妥，一字不隐，遂回明王夫人。}
}
{
    \eA{[English source not present in the checked-in Bencraft/Joly files.]}
}
{
}

\Verse{46}
{
    \zA{王夫人吃 了一惊。}
}
{
    \eA{[English source not present in the checked-in Bencraft/Joly files.]}
}
{
}

\Verse{47}
{
    \zA{想到司棋系迎春丫头，乃系那边的人，只得令人去回邢氏。}
}
{
    \eA{[English source not present in the checked-in Bencraft/Joly files.]}
}
{
}

\Verse{48}
{
    \zA{周瑞家的回道： “前日那边太太嗔着王善保家的多事，打了几个嘴巴子，如今他也装病在家，不肯 出头了。}
}
{
    \eA{[English source not present in the checked-in Bencraft/Joly files.]}
}
{
}

\Verse{49}
{
    \zA{况且又是他外孙女儿，自己打了嘴，他只好装个忘了，日久平服了再说。}
}
{
    \eA{[English source not present in the checked-in Bencraft/Joly files.]}
}
{
}

\Verse{50}
{
    \zA{如今我们过去回时，恐怕又多心，倒像咱们多事是的。}
}
{
    \eA{[English source not present in the checked-in Bencraft/Joly files.]}
}
{
}

\Verse{51}
{
    \zA{不如直把司棋带过去，一并 连赃证与那边太太瞧了，不过打一顿配了人，再指个丫头来，岂不省事？}
}
{
    \eA{[English source not present in the checked-in Bencraft/Joly files.]}
}
{
}

\Verse{52}
{
    \zA{如今白告 诉去，那边太太再推三阻四的，又说：‘既这样，你太太就该料理，又来说什么呢？’}
}
{
    \eA{[English source not present in the checked-in Bencraft/Joly files.]}
}
{
}

\Verse{53}
{
    \zA{岂不倒耽搁了？}
}
{
    \eA{[English source not present in the checked-in Bencraft/Joly files.]}
}
{
}

\Verse{54}
{
    \zA{倘或那丫头瞅空儿寻了死，反不好了。}
}
{
    \eA{[English source not present in the checked-in Bencraft/Joly files.]}
}
{
}

\Verse{55}
{
    \zA{如今看了两三天，都有些偷 懒，倘一时不到，岂不倒弄出事来？”}
}
{
    \eA{[English source not present in the checked-in Bencraft/Joly files.]}
}
{
}

\Verse{56}
{
    \zA{王夫人想了一想，说：“这也倒是。}
}
{
    \eA{[English source not present in the checked-in Bencraft/Joly files.]}
}
{
}

\Verse{57}
{
    \zA{快办了 这一件，再办咱们家的那些妖精。”}
}
{
    \eA{[English source not present in the checked-in Bencraft/Joly files.]}
}
{
}

\Verse{58}
{
    \zA{周瑞家的听说，会齐了那边几个媳妇，先到迎春房里，回明迎春。}
}
{
    \eA{[English source not present in the checked-in Bencraft/Joly files.]}
}
{
}

\Verse{59}
{
    \zA{迎春听了， 含泪似有不舍之意，因前夜之事，丫头们悄悄说了原故，虽数年之情难舍，但事关 风化，亦无可如何了。}
}
{
    \eA{[English source not present in the checked-in Bencraft/Joly files.]}
}
{
}

\Verse{60}
{
    \zA{那司棋也曾求了迎春，实指望能救，只是迎春语言迟慢，耳 软心活，是不能作主的。}
}
{
    \eA{[English source not present in the checked-in Bencraft/Joly files.]}
}
{
}

\Verse{61}
{
    \zA{司棋见了这般，知不能免，因跪着哭道：“姑娘好狠心! 哄了我这两日，如今怎么连一句话也没有？”}
}
{
    \eA{[English source not present in the checked-in Bencraft/Joly files.]}
}
{
}

\Verse{62}
{
    \zA{周瑞家的说道：“你还要姑娘留你不 成？}
}
{
    \eA{[English source not present in the checked-in Bencraft/Joly files.]}
}
{
}

\Verse{63}
{
    \zA{便留下，你也难见园里的人了。}
}
{
    \eA{[English source not present in the checked-in Bencraft/Joly files.]}
}
{
}

\Verse{64}
{
    \zA{依我们的好话，快快收了这样子，倒是人不知 鬼不觉的去罢，大家体面些。”}
}
{
    \eA{[English source not present in the checked-in Bencraft/Joly files.]}
}
{
}

\Verse{65}
{
    \zA{迎春手里拿着一本书正看呢，听了这话，书也不看， 话也不答，只管扭着身子呆呆的坐着。}
}
{
    \eA{[English source not present in the checked-in Bencraft/Joly files.]}
}
{
}

\Verse{66}
{
    \zA{周瑞家的又催道：“这么大女孩儿，自己作 的还不知道？}
}
{
    \eA{[English source not present in the checked-in Bencraft/Joly files.]}
}
{
}

\Verse{67}
{
    \zA{把姑娘都带的不好了，你还敢紧着缠磨他！”}
}
{
    \eA{[English source not present in the checked-in Bencraft/Joly files.]}
}
{
}

\Verse{68}
{
    \zA{迎春听了，方发话道： “你瞧入画也是几年的，怎么说去就去了？}
}
{
    \eA{[English source not present in the checked-in Bencraft/Joly files.]}
}
{
}

\Verse{69}
{
    \zA{自然不止你两个，想这园里凡大的都要 去呢。}
}
{
    \eA{[English source not present in the checked-in Bencraft/Joly files.]}
}
{
}

\Verse{70}
{
    \zA{依我说，将来总有一散，不如各人去罢。”}
}
{
    \eA{[English source not present in the checked-in Bencraft/Joly files.]}
}
{
}

\Verse{71}
{
    \zA{周瑞家的道：“所以到底是姑娘 明白。}
}
{
    \eA{[English source not present in the checked-in Bencraft/Joly files.]}
}
{
}

\Verse{72}
{
    \zA{明儿还有打发的人呢，你放心罢。”}
}
{
    \eA{[English source not present in the checked-in Bencraft/Joly files.]}
}
{
}

\Verse{73}
{
    \zA{司棋无法，只得含泪给迎春磕头，和众 人告别。}
}
{
    \eA{[English source not present in the checked-in Bencraft/Joly files.]}
}
{
}

\Verse{74}
{
    \zA{又向迎春耳边说：“好歹打听我受罪，替我说个情儿，就是主仆一场！”}
}
{
    \eA{[English source not present in the checked-in Bencraft/Joly files.]}
}
{
}

\Verse{75}
{
    \zA{迎春亦含泪答应：“放心。”}
}
{
    \eA{[English source not present in the checked-in Bencraft/Joly files.]}
}
{
}

\Verse{76}
{
    \zA{于是周瑞家的等人带了司棋出去，又有两个婆子将司棋所有的东西都与他拿 着。}
}
{
    \eA{[English source not present in the checked-in Bencraft/Joly files.]}
}
{
}

\Verse{77}
{
    \zA{走了没几步，只见后头绣橘赶来，一面也擦着泪，一面递给司棋一个绢包，说： “这是姑娘给你的。}
}
{
    \eA{[English source not present in the checked-in Bencraft/Joly files.]}
}
{
}

\Verse{78}
{
    \zA{主仆一场，如今一旦分离，这个给你做个念心儿罢。”}
}
{
    \eA{[English source not present in the checked-in Bencraft/Joly files.]}
}
{
}

\Verse{79}
{
    \zA{司棋接 了，不觉更哭起来了，又和绣橘哭了一回。}
}
{
    \eA{[English source not present in the checked-in Bencraft/Joly files.]}
}
{
}

\Verse{80}
{
    \zA{周瑞家的不耐烦，只管催促，二人只得 散了。}
}
{
    \eA{[English source not present in the checked-in Bencraft/Joly files.]}
}
{
}

\Verse{81}
{
    \zA{司棋因又哭告道：“婶子大娘们，好歹略徇个情儿：如今且歇一歇，让我到 相好姊妹跟前辞一辞，也是这几年我们相好一场。”}
}
{
    \eA{[English source not present in the checked-in Bencraft/Joly files.]}
}
{
}
\ChapterImage{img/chapters/147第77回 連？證攆出大觀園.jpg}


\Verse{82}
{
    \zA{周瑞家的等人皆各有事，做这 些事便是不得已了，况且又深恨他们素日大样，如今那里工夫听他的话？}
}
{
    \eA{[English source not present in the checked-in Bencraft/Joly files.]}
}
{
}

\Verse{83}
{
    \zA{因冷笑道： “我劝你去罢，别拉拉扯扯的了，我们还有正经事呢。}
}
{
    \eA{[English source not present in the checked-in Bencraft/Joly files.]}
}
{
}

\Verse{84}
{
    \zA{谁是你一个衣胞里爬出来的？}
}
{
    \eA{[English source not present in the checked-in Bencraft/Joly files.]}
}
{
}

\Verse{85}
{
    \zA{辞他们做什么？}
}
{
    \eA{[English source not present in the checked-in Bencraft/Joly files.]}
}
{
}

\Verse{86}
{
    \zA{你不过挨一会是一会，难道算了不成？}
}
{
    \eA{[English source not present in the checked-in Bencraft/Joly files.]}
}
{
}

\Verse{87}
{
    \zA{依我说，快去罢！”}
}
{
    \eA{[English source not present in the checked-in Bencraft/Joly files.]}
}
{
}

\Verse{88}
{
    \zA{一面说， 一面总不住脚，直带着出后角门去。}
}
{
    \eA{[English source not present in the checked-in Bencraft/Joly files.]}
}
{
}

\Verse{89}
{
    \zA{司棋无奈，又不敢再说，只得跟着出来。}
}
{
    \eA{[English source not present in the checked-in Bencraft/Joly files.]}
}
{
}

\Verse{90}
{
    \zA{可巧正值宝玉从外头进来，一见带了司棋出去，又见后面抱着许多东西，料着 此去再不能来了。}
}
{
    \eA{[English source not present in the checked-in Bencraft/Joly files.]}
}
{
}

\Verse{91}
{
    \zA{因听见上夜的事，并晴雯的病也因那日加重，细问晴雯，又不说 是为何。}
}
{
    \eA{[English source not present in the checked-in Bencraft/Joly files.]}
}
{
}

\Verse{92}
{
    \zA{今见司棋亦走，不觉如丧魂魄，因忙拦住问道：“那里去？”}
}
{
    \eA{[English source not present in the checked-in Bencraft/Joly files.]}
}
{
}

\Verse{93}
{
    \zA{周瑞家的等 皆知宝玉素昔行为，又恐唠叨误事，因笑道：“不干你事，快念书去罢。”}
}
{
    \eA{[English source not present in the checked-in Bencraft/Joly files.]}
}
{
}

\Verse{94}
{
    \zA{宝玉笑 道：“姐姐们且站一站，我有道理。”}
}
{
    \eA{[English source not present in the checked-in Bencraft/Joly files.]}
}
{
}

\Verse{95}
{
    \zA{周瑞家的便道：“太太吩咐不许少捱时刻。}
}
{
    \eA{[English source not present in the checked-in Bencraft/Joly files.]}
}
{
}

\Verse{96}
{
    \zA{又有什么道理？}
}
{
    \eA{[English source not present in the checked-in Bencraft/Joly files.]}
}
{
}

\Verse{97}
{
    \zA{我们只知道太太的话，管不得许多。”}
}
{
    \eA{[English source not present in the checked-in Bencraft/Joly files.]}
}
{
}

\Verse{98}
{
    \zA{司棋见了宝玉，因拉住哭道： “他们做不得主，好歹求求太太去！”}
}
{
    \eA{[English source not present in the checked-in Bencraft/Joly files.]}
}
{
}

\Verse{99}
{
    \zA{宝玉不禁也伤心，含泪说道：“我不知你做 了什么大事!晴雯也气病着，如今你又要去了，这却怎么着好！”}
}
{
    \eA{[English source not present in the checked-in Bencraft/Joly files.]}
}
{
}

\Verse{100}
{
    \zA{周瑞家的发躁向 司棋道：“你如今不是副小姐了，要不听说，我就打得你了。}
}
{
    \eA{[English source not present in the checked-in Bencraft/Joly files.]}
}
{
}

\Verse{101}
{
    \zA{别想往日有姑娘护着， 任你们作耗!越说着，还不好生走。}
}
{
    \eA{[English source not present in the checked-in Bencraft/Joly files.]}
}
{
}

\Verse{102}
{
    \zA{一个小爷见了面，也拉拉扯扯的，什么意思！”}
}
{
    \eA{[English source not present in the checked-in Bencraft/Joly files.]}
}
{
}

\Verse{103}
{
    \zA{那几个妇人不由分说，拉着司棋，便出去了。}
}
{
    \eA{[English source not present in the checked-in Bencraft/Joly files.]}
}
{
}

\Verse{104}
{
    \zA{宝玉又恐他们去告舌，恨的只瞪着他们。}
}
{
    \eA{[English source not present in the checked-in Bencraft/Joly files.]}
}
{
}

\Verse{105}
{
    \zA{看走远了，方指着恨道：“奇怪，奇 怪!怎么这些人只一嫁了汉子，染了男人的气味，就这样混账起来，比男人更可杀 了！”}
}
{
    \eA{[English source not present in the checked-in Bencraft/Joly files.]}
}
{
}

\Verse{106}
{
    \zA{守园门的婆子听了，也不禁好笑起来，因问道：“这样说，凡女儿个个是好 的了，女人个个是坏的了？”}
}
{
    \eA{[English source not present in the checked-in Bencraft/Joly files.]}
}
{
}

\Verse{107}
{
    \zA{宝玉发恨道：“不错，不错！”}
}
{
    \eA{[English source not present in the checked-in Bencraft/Joly files.]}
}
{
}

\Verse{108}
{
    \zA{正说着，只见几个老 婆子走来，忙说道：“你们小心传齐了伺候着。}
}
{
    \eA{[English source not present in the checked-in Bencraft/Joly files.]}
}
{
}

\Verse{109}
{
    \zA{此刻太太亲自到园里查人呢。”}
}
{
    \eA{[English source not present in the checked-in Bencraft/Joly files.]}
}
{
}

\Verse{110}
{
    \zA{又 吩咐：“快叫怡红院晴雯姑娘的哥嫂来，在这里等着，领出他妹子去。”}
}
{
    \eA{[English source not present in the checked-in Bencraft/Joly files.]}
}
{
}

\Verse{111}
{
    \zA{因又笑道： “阿弥陀佛!今日天睁了眼，把这个祸害妖精退送了，大家清净些。”}
}
{
    \eA{[English source not present in the checked-in Bencraft/Joly files.]}
}
{
}

\Verse{112}
{
    \zA{宝玉一闻得 王夫人进来亲查，便料道晴雯也保不住了，早飞也似的赶了去，所以后来趁愿之话， 竟未听见。}
}
{
    \eA{[English source not present in the checked-in Bencraft/Joly files.]}
}
{
}

\Verse{113}
{
    \zA{宝玉及到了怡红院，只见一群人在那里。}
}
{
    \eA{[English source not present in the checked-in Bencraft/Joly files.]}
}
{
}

\Verse{114}
{
    \zA{王夫人在屋里坐着，一脸怒色，见宝 玉也不理。}
}
{
    \eA{[English source not present in the checked-in Bencraft/Joly files.]}
}
{
}

\Verse{115}
{
    \zA{晴雯四五日水米不曾沾牙，如今现打炕上拉下来，蓬头垢面的，两个女 人搀架起来去了。}
}
{
    \eA{[English source not present in the checked-in Bencraft/Joly files.]}
}
{
}

\Verse{116}
{
    \zA{王夫人吩咐：“把他贴身的衣服撂出去，馀者留下，给好的丫头 们穿。”}
}
{
    \eA{[English source not present in the checked-in Bencraft/Joly files.]}
}
{
}

\Verse{117}
{
    \zA{又命：“把这里所有的丫头们都叫来！”}
}
{
    \eA{[English source not present in the checked-in Bencraft/Joly files.]}
}
{
}

\Verse{118}
{
    \zA{一一过目。}
}
{
    \eA{[English source not present in the checked-in Bencraft/Joly files.]}
}
{
}

\Verse{119}
{
    \zA{原来王夫人惟怕丫头们教坏了宝玉，乃从袭人起以至于极小的粗活小丫头们， 个个亲自看了一遍。}
}
{
    \eA{[English source not present in the checked-in Bencraft/Joly files.]}
}
{
}

\Verse{120}
{
    \zA{因问：“谁是和宝玉一日的生日？”}
}
{
    \eA{[English source not present in the checked-in Bencraft/Joly files.]}
}
{
}

\Verse{121}
{
    \zA{本人不敢答言。}
}
{
    \eA{[English source not present in the checked-in Bencraft/Joly files.]}
}
{
}

\Verse{122}
{
    \zA{李嬷嬷指 道：“这一个蕙香，又叫做四儿的，是同宝玉一日生日的。”}
}
{
    \eA{[English source not present in the checked-in Bencraft/Joly files.]}
}
{
}

\Verse{123}
{
    \zA{王夫人细看了一看， 虽比不上晴雯一半，却有几分水秀，视其行止，聪明皆露在外面，且也打扮的不同。}
}
{
    \eA{[English source not present in the checked-in Bencraft/Joly files.]}
}
{
}

\Verse{124}
{
    \zA{王夫人冷笑道：“这也是个没廉耻的货!他背地里说的同日生日就是夫妻，这可是 你说的？}
}
{
    \eA{[English source not present in the checked-in Bencraft/Joly files.]}
}
{
}

\Verse{125}
{
    \zA{打量我隔的远，都不知道呢。}
}
{
    \eA{[English source not present in the checked-in Bencraft/Joly files.]}
}
{
}

\Verse{126}
{
    \zA{可知我身子虽不大来，我的心耳神意时时都 在这里。}
}
{
    \eA{[English source not present in the checked-in Bencraft/Joly files.]}
}
{
}

\Verse{127}
{
    \zA{难道我统共一个宝玉，就白放心凭你们勾引坏了不成？”}
}
{
    \eA{[English source not present in the checked-in Bencraft/Joly files.]}
}
{
}

\Verse{128}
{
    \zA{这个四儿见王夫 人说着他素日和宝玉的私语，不禁红了脸，低头垂泪。}
}
{
    \eA{[English source not present in the checked-in Bencraft/Joly files.]}
}
{
}

\Verse{129}
{
    \zA{王夫人即命：“也快把他家 人叫来，领出去配人。”}
}
{
    \eA{[English source not present in the checked-in Bencraft/Joly files.]}
}
{
}

\Verse{130}
{
    \zA{又问：“那芳官呢？”}
}
{
    \eA{[English source not present in the checked-in Bencraft/Joly files.]}
}
{
}

\Verse{131}
{
    \zA{芳官只得过来。}
}
{
    \eA{[English source not present in the checked-in Bencraft/Joly files.]}
}
{
}

\Verse{132}
{
    \zA{王夫人道：“唱戏 的女孩子，自然更是狐狸精了!上次放你们，你们又不愿去，可就该安分守己才是。}
}
{
    \eA{[English source not present in the checked-in Bencraft/Joly files.]}
}
{
}

\Verse{133}
{
    \zA{你就成精鼓捣起来，调唆宝玉，无所不为！”}
}
{
    \eA{[English source not present in the checked-in Bencraft/Joly files.]}
}
{
}

\Verse{134}
{
    \zA{芳官等辩道：“并不敢调唆什么了。”}
}
{
    \eA{[English source not present in the checked-in Bencraft/Joly files.]}
}
{
}

\Verse{135}
{
    \zA{王夫人笑道：“你还强嘴!你连你干娘都压倒了，岂止别人。”}
}
{
    \eA{[English source not present in the checked-in Bencraft/Joly files.]}
}
{
}

\Verse{136}
{
    \zA{因喝命：“唤他干 娘来领去!就赏他外头找个女婿罢。}
}
{
    \eA{[English source not present in the checked-in Bencraft/Joly files.]}
}
{
}

\Verse{137}
{
    \zA{他的东西，一概给他。”}
}
{
    \eA{[English source not present in the checked-in Bencraft/Joly files.]}
}
{
}

\Verse{138}
{
    \zA{吩咐：“上年凡有姑 娘分的唱戏女孩子们，一概不许留在园里，都令其各人干娘带出，自行聘嫁。”}
}
{
    \eA{[English source not present in the checked-in Bencraft/Joly files.]}
}
{
}

\Verse{139}
{
    \zA{一 语传出，这些干娘皆感恩趁愿不尽，都约齐给王夫人磕头领去。}
}
{
    \eA{[English source not present in the checked-in Bencraft/Joly files.]}
}
{
}

\Verse{140}
{
    \zA{王夫人又满屋里搜检宝玉之物。}
}
{
    \eA{[English source not present in the checked-in Bencraft/Joly files.]}
}
{
}

\Verse{141}
{
    \zA{凡略有眼生之物，一并命收卷起来，拿到自己 房里去了。}
}
{
    \eA{[English source not present in the checked-in Bencraft/Joly files.]}
}
{
}

\Verse{142}
{
    \zA{因说：“这才干净，省得旁人口舌。”}
}
{
    \eA{[English source not present in the checked-in Bencraft/Joly files.]}
}
{
}

\Verse{143}
{
    \zA{又吩咐袭人麝月等人：“你们小 心，往后再有一点分外之事，我一概不饶!因叫人查看了，今年不宜迁挪，暂且挨 过今年，明年一并给我仍旧搬出去，才心净。”}
}
{
    \eA{[English source not present in the checked-in Bencraft/Joly files.]}
}
{
}

\Verse{144}
{
    \zA{说毕，茶也不吃，遂带领众人，又 往别处去阅人。}
}
{
    \eA{[English source not present in the checked-in Bencraft/Joly files.]}
}
{
}

\Verse{145}
{
    \zA{暂且说不到后文，如今且说宝玉只道王夫人不过来搜检搜检，无甚大事，谁知 竟这样雷嗔电怒的来了。}
}
{
    \eA{[English source not present in the checked-in Bencraft/Joly files.]}
}
{
}

\Verse{146}
{
    \zA{所责之事，皆系平日私语，一字不爽，料必不能挽回的。}
}
{
    \eA{[English source not present in the checked-in Bencraft/Joly files.]}
}
{
}

\Verse{147}
{
    \zA{虽心下恨不能一死，但王夫人盛怒之际，自不敢多言。}
}
{
    \eA{[English source not present in the checked-in Bencraft/Joly files.]}
}
{
}

\Verse{148}
{
    \zA{一直跟送王夫人到沁芳亭， 王夫人命：“回去好生念念那书!仔细明儿问你。}
}
{
    \eA{[English source not present in the checked-in Bencraft/Joly files.]}
}
{
}

\Verse{149}
{
    \zA{才已发下狠了。”}
}
{
    \eA{[English source not present in the checked-in Bencraft/Joly files.]}
}
{
}

\Verse{150}
{
    \zA{宝玉听如此说， 才回来。}
}
{
    \eA{[English source not present in the checked-in Bencraft/Joly files.]}
}
{
}

\Verse{151}
{
    \zA{一路打算：“谁这样犯舌？}
}
{
    \eA{[English source not present in the checked-in Bencraft/Joly files.]}
}
{
}

\Verse{152}
{
    \zA{况这里事也无人知道，如何就都说着了？”}
}
{
    \eA{[English source not present in the checked-in Bencraft/Joly files.]}
}
{
}

\Verse{153}
{
    \zA{一 面想，一面进来，只见袭人在那里垂泪，且去了第一等的人，岂不伤心？}
}
{
    \eA{[English source not present in the checked-in Bencraft/Joly files.]}
}
{
}

\Verse{154}
{
    \zA{便倒在床 上大哭起来。}
}
{
    \eA{[English source not present in the checked-in Bencraft/Joly files.]}
}
{
}

\Verse{155}
{
    \zA{袭人知他心里别的犹可，独有晴雯是第一件大事，乃劝道：“哭也不中用。}
}
{
    \eA{[English source not present in the checked-in Bencraft/Joly files.]}
}
{
}

\Verse{156}
{
    \zA{你 起来，我告诉你：晴雯已经好了，他这一家去，倒心净养几天。}
}
{
    \eA{[English source not present in the checked-in Bencraft/Joly files.]}
}
{
}

\Verse{157}
{
    \zA{你果然舍不得他， 等太太气消了，你再求老太太，慢慢的叫进来，也不难。}
}
{
    \eA{[English source not present in the checked-in Bencraft/Joly files.]}
}
{
}

\Verse{158}
{
    \zA{太太不过偶然听了别人的 闲言，在气头上罢了。”}
}
{
    \eA{[English source not present in the checked-in Bencraft/Joly files.]}
}
{
}

\Verse{159}
{
    \zA{宝玉道：“我究竟不知晴雯犯了什么迷天大罪！”}
}
{
    \eA{[English source not present in the checked-in Bencraft/Joly files.]}
}
{
}

\Verse{160}
{
    \zA{袭人道： “太太只嫌他生的太好了，未免轻狂些。}
}
{
    \eA{[English source not present in the checked-in Bencraft/Joly files.]}
}
{
}

\Verse{161}
{
    \zA{太太是深知这样美人似的人，心里是不能 安静的，所以很嫌他。}
}
{
    \eA{[English source not present in the checked-in Bencraft/Joly files.]}
}
{
}

\Verse{162}
{
    \zA{像我们这粗粗笨笨的倒好。”}
}
{
    \eA{[English source not present in the checked-in Bencraft/Joly files.]}
}
{
}
\ChapterImage{img/chapters/148第77回 俏丫鬟抱屈夭風流.jpg}


\Verse{163}
{
    \zA{宝玉道：“美人似的，心里就 不安静么？}
}
{
    \eA{[English source not present in the checked-in Bencraft/Joly files.]}
}
{
}

\Verse{164}
{
    \zA{你那里知道，古来美人安静的多着呢。}
}
{
    \eA{[English source not present in the checked-in Bencraft/Joly files.]}
}
{
}

\Verse{165}
{
    \zA{这也罢了，咱们私自玩话，怎么 也知道了？}
}
{
    \eA{[English source not present in the checked-in Bencraft/Joly files.]}
}
{
}

\Verse{166}
{
    \zA{又没外人走风，这可奇怪了。”}
}
{
    \eA{[English source not present in the checked-in Bencraft/Joly files.]}
}
{
}

\Verse{167}
{
    \zA{袭人道：“你有什么忌讳的？}
}
{
    \eA{[English source not present in the checked-in Bencraft/Joly files.]}
}
{
}

\Verse{168}
{
    \zA{一时高兴， 你就不管有人没人了。}
}
{
    \eA{[English source not present in the checked-in Bencraft/Joly files.]}
}
{
}

\Verse{169}
{
    \zA{我也曾使过眼色，也曾递过暗号，被那人知道了，你还不觉。”}
}
{
    \eA{[English source not present in the checked-in Bencraft/Joly files.]}
}
{
}

\Verse{170}
{
    \zA{宝玉道：“怎么人人的不是，太太都知道了，单不挑出你和麝月秋纹来？”}
}
{
    \eA{[English source not present in the checked-in Bencraft/Joly files.]}
}
{
}

\Verse{171}
{
    \zA{袭人听 了这话，心内一动，低头半日，无可回答，因便笑道：“正是呢。}
}
{
    \eA{[English source not present in the checked-in Bencraft/Joly files.]}
}
{
}

\Verse{172}
{
    \zA{若论我们，也有 玩笑不留心的去处，怎么太太竟忘了？}
}
{
    \eA{[English source not present in the checked-in Bencraft/Joly files.]}
}
{
}

\Verse{173}
{
    \zA{想是还有别的事，等完了再发放我们也未可 知。”}
}
{
    \eA{[English source not present in the checked-in Bencraft/Joly files.]}
}
{
}

\Verse{174}
{
    \zA{宝玉笑道：“你是头一个出了名的至善至贤的人，他两个又是你陶冶教育的， 焉得有什么该罚之处？}
}
{
    \eA{[English source not present in the checked-in Bencraft/Joly files.]}
}
{
}

\Verse{175}
{
    \zA{只是芳官尚小，过于伶俐些，未免倚强压倒了人，惹人厌。}
}
{
    \eA{[English source not present in the checked-in Bencraft/Joly files.]}
}
{
}

\Verse{176}
{
    \zA{四儿是我误了他：还是那年我和你拌嘴的那日起，叫上来做细活的。}
}
{
    \eA{[English source not present in the checked-in Bencraft/Joly files.]}
}
{
}

\Verse{177}
{
    \zA{众人见我待他 好，未免夺了地位，也是有的，故有今日。}
}
{
    \eA{[English source not present in the checked-in Bencraft/Joly files.]}
}
{
}

\Verse{178}
{
    \zA{只是晴雯，也是和你们一样从小儿在老 太太屋里过来的，虽生的比人强些，也没什么妨碍着谁的去处。}
}
{
    \eA{[English source not present in the checked-in Bencraft/Joly files.]}
}
{
}

\Verse{179}
{
    \zA{就只是他的性情爽 利，口角锋芒，竟也没见他得罪了那一个。}
}
{
    \eA{[English source not present in the checked-in Bencraft/Joly files.]}
}
{
}

\Verse{180}
{
    \zA{可是你说的，想是他过于生得好了，反 被这个好带累了！”}
}
{
    \eA{[English source not present in the checked-in Bencraft/Joly files.]}
}
{
}

\Verse{181}
{
    \zA{说毕，复又哭起来。}
}
{
    \eA{[English source not present in the checked-in Bencraft/Joly files.]}
}
{
}

\Verse{182}
{
    \zA{袭人细揣，此话只是宝玉有疑他之意，竟不好再劝，因叹道：“天知道罢了。}
}
{
    \eA{[English source not present in the checked-in Bencraft/Joly files.]}
}
{
}

\Verse{183}
{
    \zA{此时也查不出人来了。}
}
{
    \eA{[English source not present in the checked-in Bencraft/Joly files.]}
}
{
}

\Verse{184}
{
    \zA{白哭一会子，也无益了。”}
}
{
    \eA{[English source not present in the checked-in Bencraft/Joly files.]}
}
{
}

\Verse{185}
{
    \zA{宝玉冷笑道：“原是想他自幼娇 生惯养的，何尝受过一日委屈？}
}
{
    \eA{[English source not present in the checked-in Bencraft/Joly files.]}
}
{
}

\Verse{186}
{
    \zA{如今是一盆才透出嫩箭的兰花送到猪圈里去一般。}
}
{
    \eA{[English source not present in the checked-in Bencraft/Joly files.]}
}
{
}

\Verse{187}
{
    \zA{况又是一身重病，里头一肚子闷气。}
}
{
    \eA{[English source not present in the checked-in Bencraft/Joly files.]}
}
{
}

\Verse{188}
{
    \zA{他又没有亲爹热娘，只有一个醉泥鳅姑舅哥哥， 他这一去，那里还等得一月半月？}
}
{
    \eA{[English source not present in the checked-in Bencraft/Joly files.]}
}
{
}

\Verse{189}
{
    \zA{再不能见一面两面的了！”}
}
{
    \eA{[English source not present in the checked-in Bencraft/Joly files.]}
}
{
}

\Verse{190}
{
    \zA{说着，越发心痛起来。}
}
{
    \eA{[English source not present in the checked-in Bencraft/Joly files.]}
}
{
}

\Verse{191}
{
    \zA{袭人笑道：“可是你‘自许州官放火，不许百姓点灯’。}
}
{
    \eA{[English source not present in the checked-in Bencraft/Joly files.]}
}
{
}

\Verse{192}
{
    \zA{我们偶说一句妨碍的话， 你就说不吉利；你如今好好的咒他，就该的了？”}
}
{
    \eA{[English source not present in the checked-in Bencraft/Joly files.]}
}
{
}

\Verse{193}
{
    \zA{宝玉道：“我不是妄口咒人，今 年春天已有兆头的。”}
}
{
    \eA{[English source not present in the checked-in Bencraft/Joly files.]}
}
{
}

\Verse{194}
{
    \zA{袭人忙问：“何兆？”}
}
{
    \eA{[English source not present in the checked-in Bencraft/Joly files.]}
}
{
}

\Verse{195}
{
    \zA{宝玉道：“这阶下好好的一株海棠花， 竟无故死了半边，我就知道有坏事，果然应在他身上。”}
}
{
    \eA{[English source not present in the checked-in Bencraft/Joly files.]}
}
{
}

\Verse{196}
{
    \zA{袭人听了，又笑起来说： “我要不说，又掌不住，你也太婆婆妈妈的了。}
}
{
    \eA{[English source not present in the checked-in Bencraft/Joly files.]}
}
{
}

\Verse{197}
{
    \zA{这样的话，怎么是你读书的人说 的？”}
}
{
    \eA{[English source not present in the checked-in Bencraft/Joly files.]}
}
{
}

\Verse{198}
{
    \zA{宝玉叹道：“你们那里知道？}
}
{
    \eA{[English source not present in the checked-in Bencraft/Joly files.]}
}
{
}

\Verse{199}
{
    \zA{不但草木，凡天下有情有理的东西，也和人一 样，得了知己，便极有灵验的。}
}
{
    \eA{[English source not present in the checked-in Bencraft/Joly files.]}
}
{
}

\Verse{200}
{
    \zA{若用大题目比，就像孔子庙前桧树，坟前的蓍草， 诸葛祠前的柏树，岳武穆坟前的松树：这都是堂堂正大之气，千古不磨之物。}
}
{
    \eA{[English source not present in the checked-in Bencraft/Joly files.]}
}
{
}

\Verse{201}
{
    \zA{世乱 他就枯干了，世治他就茂盛了，凡千年枯了又生的几次，这不是应兆么？}
}
{
    \eA{[English source not present in the checked-in Bencraft/Joly files.]}
}
{
}

\Verse{202}
{
    \zA{若是小题 目比，就像杨太真沈香亭的木芍药，端正楼的相思树，王昭君坟上的长青草，难道 不也有灵验？}
}
{
    \eA{[English source not present in the checked-in Bencraft/Joly files.]}
}
{
}

\Verse{203}
{
    \zA{所以这海棠亦是应着人生的。”}
}
{
    \eA{[English source not present in the checked-in Bencraft/Joly files.]}
}
{
}

\Verse{204}
{
    \zA{袭人听了这篇痴话，又可笑，又可叹， 因笑道：“真真的这话越发说上我的气来了。}
}
{
    \eA{[English source not present in the checked-in Bencraft/Joly files.]}
}
{
}

\Verse{205}
{
    \zA{那晴雯是个什么东西？}
}
{
    \eA{[English source not present in the checked-in Bencraft/Joly files.]}
}
{
}

\Verse{206}
{
    \zA{就费这样心思， 比出这些正经人来。}
}
{
    \eA{[English source not present in the checked-in Bencraft/Joly files.]}
}
{
}

\Verse{207}
{
    \zA{还有一说：他纵好，也越不过我的次序去。}
}
{
    \eA{[English source not present in the checked-in Bencraft/Joly files.]}
}
{
}

\Verse{208}
{
    \zA{就是这海棠，也该 先来比我，也还轮不到他。}
}
{
    \eA{[English source not present in the checked-in Bencraft/Joly files.]}
}
{
}

\Verse{209}
{
    \zA{想是我要死的了。”}
}
{
    \eA{[English source not present in the checked-in Bencraft/Joly files.]}
}
{
}

\Verse{210}
{
    \zA{宝玉听说，忙掩他的嘴，劝道：“这是何苦？}
}
{
    \eA{[English source not present in the checked-in Bencraft/Joly files.]}
}
{
}

\Verse{211}
{
    \zA{一个未是，你又这样起来。}
}
{
    \eA{[English source not present in the checked-in Bencraft/Joly files.]}
}
{
}

\Verse{212}
{
    \zA{罢了， 再别提这事，别弄的去了三个，又饶上一个。”}
}
{
    \eA{[English source not present in the checked-in Bencraft/Joly files.]}
}
{
}

\Verse{213}
{
    \zA{袭人听说，心下暗喜道：“若不如 此，也没个了局。”}
}
{
    \eA{[English source not present in the checked-in Bencraft/Joly files.]}
}
{
}

\Verse{214}
{
    \zA{宝玉又道：“我还有一句话要和你商量，不知你肯不肯：现在 他的东西，是瞒上不瞒下，悄悄的送还他去。}
}
{
    \eA{[English source not present in the checked-in Bencraft/Joly files.]}
}
{
}

\Verse{215}
{
    \zA{再或有咱们常日积攒下的钱，拿几吊 出去，给他养病，也是你姐妹好了一场。”}
}
{
    \eA{[English source not present in the checked-in Bencraft/Joly files.]}
}
{
}

\Verse{216}
{
    \zA{袭人听了，笑道：“你太把我看得忒小 器又没人心了。}
}
{
    \eA{[English source not present in the checked-in Bencraft/Joly files.]}
}
{
}

\Verse{217}
{
    \zA{这话还等你说？}
}
{
    \eA{[English source not present in the checked-in Bencraft/Joly files.]}
}
{
}

\Verse{218}
{
    \zA{我才把他的衣裳各物已打点下了，放在那里。}
}
{
    \eA{[English source not present in the checked-in Bencraft/Joly files.]}
}
{
}

\Verse{219}
{
    \zA{如今 白日里人多眼杂，又恐生事，且等到晚上，悄悄的叫宋妈给他拿去。}
}
{
    \eA{[English source not present in the checked-in Bencraft/Joly files.]}
}
{
}

\Verse{220}
{
    \zA{我还有攒下的 几吊钱，也给他去。”}
}
{
    \eA{[English source not present in the checked-in Bencraft/Joly files.]}
}
{
}

\Verse{221}
{
    \zA{宝玉听了，点点头儿。}
}
{
    \eA{[English source not present in the checked-in Bencraft/Joly files.]}
}
{
}

\Verse{222}
{
    \zA{袭人笑道：“我原是久已‘出名的贤 人’，连这一点子好名还不会买去不成？”}
}
{
    \eA{[English source not present in the checked-in Bencraft/Joly files.]}
}
{
}

\Verse{223}
{
    \zA{宝玉听了他方才说的，又陪笑抚慰他， 怕他寒了心。}
}
{
    \eA{[English source not present in the checked-in Bencraft/Joly files.]}
}
{
}

\Verse{224}
{
    \zA{晚间，果遣宋妈送去。}
}
{
    \eA{[English source not present in the checked-in Bencraft/Joly files.]}
}
{
}

\Verse{225}
{
    \zA{宝玉将一切人稳住，便独自得便，到园子后角门，央一个老婆子，带他到晴雯 家去。}
}
{
    \eA{[English source not present in the checked-in Bencraft/Joly files.]}
}
{
}

\Verse{226}
{
    \zA{先这婆子百般不肯，只说怕人知道，“回了太太，我还吃饭不吃饭？”}
}
{
    \eA{[English source not present in the checked-in Bencraft/Joly files.]}
}
{
}

\Verse{227}
{
    \zA{无奈 宝玉死活央告，又许他些钱，那个婆子方带了他去。}
}
{
    \eA{[English source not present in the checked-in Bencraft/Joly files.]}
}
{
}

\Verse{228}
{
    \zA{却说这晴雯当日系赖大买的。}
}
{
    \eA{[English source not present in the checked-in Bencraft/Joly files.]}
}
{
}

\Verse{229}
{
    \zA{还有个姑舅哥哥，叫做吴贵，人都叫他贵儿。}
}
{
    \eA{[English source not present in the checked-in Bencraft/Joly files.]}
}
{
}

\Verse{230}
{
    \zA{那 时晴雯才得十岁，时常赖嬷嬷带进来，贾母见了喜欢，故此赖嬷嬷就孝敬了贾母。}
}
{
    \eA{[English source not present in the checked-in Bencraft/Joly files.]}
}
{
}

\Verse{231}
{
    \zA{过了几年，赖大又给他姑舅哥哥娶了一房媳妇。}
}
{
    \eA{[English source not present in the checked-in Bencraft/Joly files.]}
}
{
}

\Verse{232}
{
    \zA{谁知贵儿一味胆小老实，那媳妇却 倒伶俐，又兼有几分姿色，看着贵儿无能为，便每日家打扮的妖妖调调，两只眼儿 水汪汪的。}
}
{
    \eA{[English source not present in the checked-in Bencraft/Joly files.]}
}
{
}

\Verse{233}
{
    \zA{招惹的赖大家人如蝇逐臭，渐渐做出些风流勾当来。}
}
{
    \eA{[English source not present in the checked-in Bencraft/Joly files.]}
}
{
}

\Verse{234}
{
    \zA{那时晴雯已在宝玉 屋里，他便央及了晴雯转求凤姐，合赖大家的要过来。}
}
{
    \eA{[English source not present in the checked-in Bencraft/Joly files.]}
}
{
}

\Verse{235}
{
    \zA{目今两口儿就在园子后角门 外居住，伺候园中买办杂差。}
}
{
    \eA{[English source not present in the checked-in Bencraft/Joly files.]}
}
{
}

\Verse{236}
{
    \zA{这晴雯一时被撵出来，住在他家。}
}
{
    \eA{[English source not present in the checked-in Bencraft/Joly files.]}
}
{
}

\Verse{237}
{
    \zA{那媳妇那里有心肠 照管？}
}
{
    \eA{[English source not present in the checked-in Bencraft/Joly files.]}
}
{
}

\Verse{238}
{
    \zA{吃了饭便自去串门子，只剩下晴雯一人，在外间屋内爬着。}
}
{
    \eA{[English source not present in the checked-in Bencraft/Joly files.]}
}
{
}

\Verse{239}
{
    \zA{宝玉命那婆子在外？}
}
{
    \eA{[English source not present in the checked-in Bencraft/Joly files.]}
}
{
}

\Verse{240}
{
    \zA{望，他独掀起布帘进来，一眼就看见晴雯睡在一领芦席 上，幸而被褥还是旧日铺盖的。}
}
{
    \eA{[English source not present in the checked-in Bencraft/Joly files.]}
}
{
}

\Verse{241}
{
    \zA{心内不知自己怎么才好，因上来含泪伸手，轻轻拉 他，悄唤两声。}
}
{
    \eA{[English source not present in the checked-in Bencraft/Joly files.]}
}
{
}

\Verse{242}
{
    \zA{当下晴雯又因着了风，又受了哥嫂的歹话，病上加病，嗽了一日， 才朦胧睡了。}
}
{
    \eA{[English source not present in the checked-in Bencraft/Joly files.]}
}
{
}
\ChapterImage{img/chapters/149第77-78回 美優伶斬情歸水月.jpg}


\Verse{243}
{
    \zA{忽闻有人唤他，强展双眸，一见是宝玉，又惊又喜，又悲又痛，一把 死攥住他的手，哽咽了半日，方说道：“我只道不得见你了！”}
}
{
    \eA{[English source not present in the checked-in Bencraft/Joly files.]}
}
{
}

\Verse{244}
{
    \zA{接着便嗽个不住。}
}
{
    \eA{[English source not present in the checked-in Bencraft/Joly files.]}
}
{
}

\Verse{245}
{
    \zA{宝玉也只有哽咽之分。}
}
{
    \eA{[English source not present in the checked-in Bencraft/Joly files.]}
}
{
}

\Verse{246}
{
    \zA{晴雯道：“阿弥陀佛，你来得好，且把那茶倒半碗我喝。}
}
{
    \eA{[English source not present in the checked-in Bencraft/Joly files.]}
}
{
}

\Verse{247}
{
    \zA{渴 了半日，叫半个人也叫不着。”}
}
{
    \eA{[English source not present in the checked-in Bencraft/Joly files.]}
}
{
}

\Verse{248}
{
    \zA{宝玉听说，忙拭泪问：“茶在那里？”}
}
{
    \eA{[English source not present in the checked-in Bencraft/Joly files.]}
}
{
}

\Verse{249}
{
    \zA{晴雯道：“在 炉台上。”}
}
{
    \eA{[English source not present in the checked-in Bencraft/Joly files.]}
}
{
}

\Verse{250}
{
    \zA{宝玉看时，虽有个黑煤乌嘴的吊子，也不像个茶壶。}
}
{
    \eA{[English source not present in the checked-in Bencraft/Joly files.]}
}
{
}

\Verse{251}
{
    \zA{只得桌上去拿一个 碗，未到手内，先闻得油膻之气。}
}
{
    \eA{[English source not present in the checked-in Bencraft/Joly files.]}
}
{
}

\Verse{252}
{
    \zA{宝玉只得拿了来，先拿些水洗了两次，复用自己 的绢子拭了，闻了闻还有些气味，没奈何，提起壶来斟了半碗。}
}
{
    \eA{[English source not present in the checked-in Bencraft/Joly files.]}
}
{
}

\Verse{253}
{
    \zA{看时绛红的也不大 像茶。}
}
{
    \eA{[English source not present in the checked-in Bencraft/Joly files.]}
}
{
}

\Verse{254}
{
    \zA{晴雯扶枕道：“快给我喝一口罢，这就是茶了。}
}
{
    \eA{[English source not present in the checked-in Bencraft/Joly files.]}
}
{
}

\Verse{255}
{
    \zA{那里比得咱们的茶呢。”}
}
{
    \eA{[English source not present in the checked-in Bencraft/Joly files.]}
}
{
}

\Verse{256}
{
    \zA{宝 玉听说，先自己尝了一尝，并无茶味，咸涩不堪，只得递给晴雯。}
}
{
    \eA{[English source not present in the checked-in Bencraft/Joly files.]}
}
{
}

\Verse{257}
{
    \zA{只见晴雯如得了 甘露一般，一气都灌下去了。}
}
{
    \eA{[English source not present in the checked-in Bencraft/Joly files.]}
}
{
}

\Verse{258}
{
    \zA{宝玉看着，眼中泪直流下来，连自己的身子都不知为何物了，一面问道：“你 有什么说的？}
}
{
    \eA{[English source not present in the checked-in Bencraft/Joly files.]}
}
{
}

\Verse{259}
{
    \zA{趁着没人，告诉我。”}
}
{
    \eA{[English source not present in the checked-in Bencraft/Joly files.]}
}
{
}

\Verse{260}
{
    \zA{晴雯呜咽道：“有什么可说的!不过是挨一刻是 一刻，挨一日是一日。}
}
{
    \eA{[English source not present in the checked-in Bencraft/Joly files.]}
}
{
}

\Verse{261}
{
    \zA{我已知横竖不过三五日的光景，我就好回去了，只是一件， 我死也不甘心：我虽生得比别人好些，并没有私情勾引你，怎么一口死咬定了我是
        个‘狐狸精’!我今儿既担了虚名，况且没了远限，不是我说一句后悔的话：早知 如此，我当日――”说到这里，气往上咽，便说不出来，两手已经冰凉。}
}
{
    \eA{[English source not present in the checked-in Bencraft/Joly files.]}
}
{
}

\Verse{262}
{
    \zA{宝玉又痛 又急，又害怕，便歪在席上，一只手攥着他的手，一只手轻轻的给他捶打着。}
}
{
    \eA{[English source not present in the checked-in Bencraft/Joly files.]}
}
{
}

\Verse{263}
{
    \zA{又不 敢大声的叫，真真万箭攒心。}
}
{
    \eA{[English source not present in the checked-in Bencraft/Joly files.]}
}
{
}

\Verse{264}
{
    \zA{两三句话时晴雯才哭出来，宝玉拉着他的手，只觉瘦 如枯柴。}
}
{
    \eA{[English source not present in the checked-in Bencraft/Joly files.]}
}
{
}

\Verse{265}
{
    \zA{腕上犹戴着四个银镯，因哭道：“除下来，等好了再戴上去罢。”}
}
{
    \eA{[English source not present in the checked-in Bencraft/Joly files.]}
}
{
}

\Verse{266}
{
    \zA{又说： “这一病好了，又伤好些！”}
}
{
    \eA{[English source not present in the checked-in Bencraft/Joly files.]}
}
{
}

\Verse{267}
{
    \zA{晴雯拭泪，把那手用力拳回，搁在口边，狠命一咬， 只听“咯吱”一声，把两根葱管一般的指甲齐根咬下，拉了宝玉的手，将指甲搁在 他手里。}
}
{
    \eA{[English source not present in the checked-in Bencraft/Joly files.]}
}
{
}

\Verse{268}
{
    \zA{又回手扎挣着，连揪带脱，在被窝内将贴身穿着的一件旧红绫小袄儿脱下， 递给宝玉。}
}
{
    \eA{[English source not present in the checked-in Bencraft/Joly files.]}
}
{
}

\Verse{269}
{
    \zA{不想虚弱透了的人，那里禁得这么抖搂，早喘成一处了。}
}
{
    \eA{[English source not present in the checked-in Bencraft/Joly files.]}
}
{
}

\Verse{270}
{
    \zA{宝玉见他这般， 已经会意，连忙解开外衣，将自己的袄儿褪下来，盖在他身上。}
}
{
    \eA{[English source not present in the checked-in Bencraft/Joly files.]}
}
{
}

\Verse{271}
{
    \zA{却把这件穿上，不 及扣钮子，只用外头衣裳掩了。}
}
{
    \eA{[English source not present in the checked-in Bencraft/Joly files.]}
}
{
}

\Verse{272}
{
    \zA{刚系腰时，只见晴雯睁眼道：“你扶起我来坐坐。”}
}
{
    \eA{[English source not present in the checked-in Bencraft/Joly files.]}
}
{
}

\Verse{273}
{
    \zA{宝玉只得扶他。}
}
{
    \eA{[English source not present in the checked-in Bencraft/Joly files.]}
}
{
}

\Verse{274}
{
    \zA{那里扶得起？}
}
{
    \eA{[English source not present in the checked-in Bencraft/Joly files.]}
}
{
}

\Verse{275}
{
    \zA{好容易欠起半身，晴雯伸手把宝玉的袄儿往自己身上 拉。}
}
{
    \eA{[English source not present in the checked-in Bencraft/Joly files.]}
}
{
}

\Verse{276}
{
    \zA{宝玉连忙给他披上，拖着？}
}
{
    \eA{[English source not present in the checked-in Bencraft/Joly files.]}
}
{
}

\Verse{277}
{
    \zA{膊，伸上袖子，轻轻放倒，然后将他的指甲装在荷 包里。}
}
{
    \eA{[English source not present in the checked-in Bencraft/Joly files.]}
}
{
}

\Verse{278}
{
    \zA{晴雯哭道：“你去罢!这里腌？}
}
{
    \eA{[English source not present in the checked-in Bencraft/Joly files.]}
}
{
}

\Verse{279}
{
    \zA{，你那里受得？}
}
{
    \eA{[English source not present in the checked-in Bencraft/Joly files.]}
}
{
}

\Verse{280}
{
    \zA{你的身子要紧。}
}
{
    \eA{[English source not present in the checked-in Bencraft/Joly files.]}
}
{
}

\Verse{281}
{
    \zA{今日这一来， 我就死了，也不枉担了虚名！”}
}
{
    \eA{[English source not present in the checked-in Bencraft/Joly files.]}
}
{
}

\Verse{282}
{
    \zA{一语未完，只见他嫂子笑嘻嘻掀帘进来道：“好呀，你两个的话，我已都听见 了。”}
}
{
    \eA{[English source not present in the checked-in Bencraft/Joly files.]}
}
{
}

\Verse{283}
{
    \zA{又向宝玉道：“你一个做主子的，跑到下人房里来做什么？}
}
{
    \eA{[English source not present in the checked-in Bencraft/Joly files.]}
}
{
}

\Verse{284}
{
    \zA{看着我年轻长的 俊，你敢只是来调戏我么？”}
}
{
    \eA{[English source not present in the checked-in Bencraft/Joly files.]}
}
{
}

\Verse{285}
{
    \zA{宝玉听见，吓得忙陪笑央及道：“好姐姐，快别大声 的。}
}
{
    \eA{[English source not present in the checked-in Bencraft/Joly files.]}
}
{
}

\Verse{286}
{
    \zA{他伏侍我一场，我私自来瞧瞧他。”}
}
{
    \eA{[English source not present in the checked-in Bencraft/Joly files.]}
}
{
}

\Verse{287}
{
    \zA{那媳妇儿点着头儿，笑道：“怨不得人家 都说你有情有义儿的。”}
}
{
    \eA{[English source not present in the checked-in Bencraft/Joly files.]}
}
{
}

\Verse{288}
{
    \zA{便一手拉了宝玉进里间来，笑道：“你要不叫我嚷，这也 容易。}
}
{
    \eA{[English source not present in the checked-in Bencraft/Joly files.]}
}
{
}

\Verse{289}
{
    \zA{你只是依我一件事。”}
}
{
    \eA{[English source not present in the checked-in Bencraft/Joly files.]}
}
{
}

\Verse{290}
{
    \zA{说着，便自己坐在炕沿上，把宝玉拉在怀中，紧紧的 将两条腿夹住。}
}
{
    \eA{[English source not present in the checked-in Bencraft/Joly files.]}
}
{
}

\Verse{291}
{
    \zA{宝玉那里见过这个？}
}
{
    \eA{[English source not present in the checked-in Bencraft/Joly files.]}
}
{
}

\Verse{292}
{
    \zA{心内早突突的跳起来了。}
}
{
    \eA{[English source not present in the checked-in Bencraft/Joly files.]}
}
{
}

\Verse{293}
{
    \zA{急得满面红胀，身上 乱战，又羞又愧又怕又恼，只说：“好姐姐，别闹。”}
}
{
    \eA{[English source not present in the checked-in Bencraft/Joly files.]}
}
{
}

\Verse{294}
{
    \zA{那媳妇乜斜了眼儿，笑道： “呸，成日家听见你在女孩儿们身上做工夫，怎么今儿个就发起讪来了？”}
}
{
    \eA{[English source not present in the checked-in Bencraft/Joly files.]}
}
{
}

\Verse{295}
{
    \zA{宝玉红 了脸，笑道：“姐姐撒开手，有话咱们慢慢儿的说。}
}
{
    \eA{[English source not present in the checked-in Bencraft/Joly files.]}
}
{
}

\Verse{296}
{
    \zA{外头有老妈妈听见，什么意思 呢？”}
}
{
    \eA{[English source not present in the checked-in Bencraft/Joly files.]}
}
{
}

\Verse{297}
{
    \zA{那媳妇那里肯放，笑道：“我早进来了，已经叫那老婆子去到园门口儿等着 呢。}
}
{
    \eA{[English source not present in the checked-in Bencraft/Joly files.]}
}
{
}

\Verse{298}
{
    \zA{我等什么儿似的，今日才等着你了!你要不依我，我就嚷起来，叫里头太太听 见了，我看你怎么样？}
}
{
    \eA{[English source not present in the checked-in Bencraft/Joly files.]}
}
{
}

\Verse{299}
{
    \zA{你这么个人，只这么大胆子儿。}
}
{
    \eA{[English source not present in the checked-in Bencraft/Joly files.]}
}
{
}

\Verse{300}
{
    \zA{我刚才进来了好一会子，在 窗下细听，屋里只你两个人，我只道有些个体己话儿。}
}
{
    \eA{[English source not present in the checked-in Bencraft/Joly files.]}
}
{
}

\Verse{301}
{
    \zA{这么看起来，你们两个人竟 还是各不相扰儿呢。}
}
{
    \eA{[English source not present in the checked-in Bencraft/Joly files.]}
}
{
}

\Verse{302}
{
    \zA{我可不能像他那么傻。”}
}
{
    \eA{[English source not present in the checked-in Bencraft/Joly files.]}
}
{
}

\Verse{303}
{
    \zA{说着，就要动手。}
}
{
    \eA{[English source not present in the checked-in Bencraft/Joly files.]}
}
{
}

\Verse{304}
{
    \zA{宝玉急的死往外拽。}
}
{
    \eA{[English source not present in the checked-in Bencraft/Joly files.]}
}
{
}

\Verse{305}
{
    \zA{正闹着，只听窗外有人问：“这晴雯姐姐在这里住呢不是？”}
}
{
    \eA{[English source not present in the checked-in Bencraft/Joly files.]}
}
{
}

\Verse{306}
{
    \zA{那媳妇子也吓了 一跳，连忙放了宝玉。}
}
{
    \eA{[English source not present in the checked-in Bencraft/Joly files.]}
}
{
}

\Verse{307}
{
    \zA{这宝玉已经吓怔了，听不出声音。}
}
{
    \eA{[English source not present in the checked-in Bencraft/Joly files.]}
}
{
}

\Verse{308}
{
    \zA{外边晴雯听见他嫂子缠磨 宝玉，又急又臊又气，一阵虚火上攻，早昏晕过去。}
}
{
    \eA{[English source not present in the checked-in Bencraft/Joly files.]}
}
{
}

\Verse{309}
{
    \zA{那媳妇连忙答应着，出来看， 不是别人，却是柳五儿和他母亲两个，抱着一个包袱。}
}
{
    \eA{[English source not present in the checked-in Bencraft/Joly files.]}
}
{
}

\Verse{310}
{
    \zA{柳家的拿着几吊钱，悄悄的 问那媳妇道：“这是里头袭姑娘叫拿出来给你们姑娘的。}
}
{
    \eA{[English source not present in the checked-in Bencraft/Joly files.]}
}
{
}

\Verse{311}
{
    \zA{他在那屋里呢？”}
}
{
    \eA{[English source not present in the checked-in Bencraft/Joly files.]}
}
{
}

\Verse{312}
{
    \zA{那媳妇 儿笑道：“就是这个屋子，那里还有屋子？”}
}
{
    \eA{[English source not present in the checked-in Bencraft/Joly files.]}
}
{
}

\Verse{313}
{
    \zA{那柳家的领着五儿刚进门来，只见一个人影儿往屋里一闪。}
}
{
    \eA{[English source not present in the checked-in Bencraft/Joly files.]}
}
{
}

\Verse{314}
{
    \zA{柳家的素知这媳妇 子不妥，只打量是他的私人。}
}
{
    \eA{[English source not present in the checked-in Bencraft/Joly files.]}
}
{
}

\Verse{315}
{
    \zA{看见晴雯睡着了，连忙放下，带着五儿便往外走。}
}
{
    \eA{[English source not present in the checked-in Bencraft/Joly files.]}
}
{
}

\Verse{316}
{
    \zA{谁 知五儿眼尖，早已见是宝玉，便问他母亲道：“头里不是袭人姐姐那里悄悄儿的找 宝二爷呢吗？”}
}
{
    \eA{[English source not present in the checked-in Bencraft/Joly files.]}
}
{
}

\Verse{317}
{
    \zA{柳家的道：“嗳哟，可是忘了。}
}
{
    \eA{[English source not present in the checked-in Bencraft/Joly files.]}
}
{
}

\Verse{318}
{
    \zA{方才老宋妈说：‘见宝二爷出角门 来了。}
}
{
    \eA{[English source not present in the checked-in Bencraft/Joly files.]}
}
{
}

\Verse{319}
{
    \zA{门上还有人等着，要关园门呢。’”}
}
{
    \eA{[English source not present in the checked-in Bencraft/Joly files.]}
}
{
}

\Verse{320}
{
    \zA{因回头问那媳妇儿。}
}
{
    \eA{[English source not present in the checked-in Bencraft/Joly files.]}
}
{
}

\Verse{321}
{
    \zA{那媳妇儿自己心虚， 便道：“宝二爷那里肯到我们这屋里来？”}
}
{
    \eA{[English source not present in the checked-in Bencraft/Joly files.]}
}
{
}

\Verse{322}
{
    \zA{柳家的听说，便要走。}
}
{
    \eA{[English source not present in the checked-in Bencraft/Joly files.]}
}
{
}

\Verse{323}
{
    \zA{这宝玉一则怕关 了门，二则怕那媳妇子进来又缠，也顾不得什么了，连忙掀了帘子出来道：“柳嫂 子，你等等我，一路儿走。”}
}
{
    \eA{[English source not present in the checked-in Bencraft/Joly files.]}
}
{
}
\ChapterImage{img/chapters/150第77-78回 老學士閑證 詞.jpg}


\Verse{324}
{
    \zA{柳家的听了，倒唬了一大跳，说：“我的爷，你怎么 跑了这里来了？”}
}
{
    \eA{[English source not present in the checked-in Bencraft/Joly files.]}
}
{
}

\Verse{325}
{
    \zA{那宝玉也不答言，一直飞走。}
}
{
    \eA{[English source not present in the checked-in Bencraft/Joly files.]}
}
{
}

\Verse{326}
{
    \zA{那五儿道：“妈妈，你快叫住宝二 爷不用忙，留神冒冒失失，被人碰见倒不好。}
}
{
    \eA{[English source not present in the checked-in Bencraft/Joly files.]}
}
{
}

\Verse{327}
{
    \zA{况且才出来时，袭人姐姐已经打发人 留了门了。”}
}
{
    \eA{[English source not present in the checked-in Bencraft/Joly files.]}
}
{
}

\Verse{328}
{
    \zA{说着，赶忙同他妈来赶宝玉。}
}
{
    \eA{[English source not present in the checked-in Bencraft/Joly files.]}
}
{
}

\Verse{329}
{
    \zA{这里晴雯的嫂子干瞅着，把个妙人儿走 了。}
}
{
    \eA{[English source not present in the checked-in Bencraft/Joly files.]}
}
{
}

\Verse{330}
{
    \zA{却说宝玉跑进角门，才把心放下来，还是突突乱跳。}
}
{
    \eA{[English source not present in the checked-in Bencraft/Joly files.]}
}
{
}

\Verse{331}
{
    \zA{又怕五儿关在外头，眼巴 巴瞅着他母女也进来了。}
}
{
    \eA{[English source not present in the checked-in Bencraft/Joly files.]}
}
{
}

\Verse{332}
{
    \zA{远远听见里边嬷嬷们正查人，若再迟一步，就关了园门了。}
}
{
    \eA{[English source not present in the checked-in Bencraft/Joly files.]}
}
{
}

\Verse{333}
{
    \zA{宝玉进入园中，且喜无人知道。}
}
{
    \eA{[English source not present in the checked-in Bencraft/Joly files.]}
}
{
}

\Verse{334}
{
    \zA{到了自己房里，告诉袭人，只说在薛姨妈家去的， 也就罢了。}
}
{
    \eA{[English source not present in the checked-in Bencraft/Joly files.]}
}
{
}

\Verse{335}
{
    \zA{一时铺床，袭人不得不问：“今日怎么睡？”}
}
{
    \eA{[English source not present in the checked-in Bencraft/Joly files.]}
}
{
}

\Verse{336}
{
    \zA{宝玉道：“不管怎么睡罢 了。”}
}
{
    \eA{[English source not present in the checked-in Bencraft/Joly files.]}
}
{
}

\Verse{337}
{
    \zA{原来这一二年来，袭人因王夫人看重了他，越发自要尊重，凡背人之处或夜 晚之间，总不与宝玉狎昵，较先小时反倒疏远了。}
}
{
    \eA{[English source not present in the checked-in Bencraft/Joly files.]}
}
{
}

\Verse{338}
{
    \zA{虽无大事办理，然一应针线，并 宝玉及诸小丫头出入银钱衣履什物等事，也甚烦琐，且有吐血之症，故近来夜间总 不与宝玉同房。}
}
{
    \eA{[English source not present in the checked-in Bencraft/Joly files.]}
}
{
}

\Verse{339}
{
    \zA{宝玉夜间胆小，醒了便要唤人，因晴雯睡卧警醒，故夜间一应茶水 起坐呼唤之事，悉皆委他一人，所以宝玉外床只是晴雯睡着。}
}
{
    \eA{[English source not present in the checked-in Bencraft/Joly files.]}
}
{
}

\Verse{340}
{
    \zA{他今去了，袭人只得 将自己铺盖搬来，铺设床外。}
}
{
    \eA{[English source not present in the checked-in Bencraft/Joly files.]}
}
{
}

\Verse{341}
{
    \zA{宝玉发了一晚上的呆。}
}
{
    \eA{[English source not present in the checked-in Bencraft/Joly files.]}
}
{
}

\Verse{342}
{
    \zA{袭人催他睡下，然后自睡。}
}
{
    \eA{[English source not present in the checked-in Bencraft/Joly files.]}
}
{
}

\Verse{343}
{
    \zA{只听宝玉在枕上长吁短叹， 覆去翻来，直至三更以后，方渐渐安顿了。}
}
{
    \eA{[English source not present in the checked-in Bencraft/Joly files.]}
}
{
}

\Verse{344}
{
    \zA{袭人方放心，也就蒙？}
}
{
    \eA{[English source not present in the checked-in Bencraft/Joly files.]}
}
{
}

\Verse{345}
{
    \zA{睡着。}
}
{
    \eA{[English source not present in the checked-in Bencraft/Joly files.]}
}
{
}

\Verse{346}
{
    \zA{没半盏茶 时，只听宝玉叫“晴雯”。}
}
{
    \eA{[English source not present in the checked-in Bencraft/Joly files.]}
}
{
}

\Verse{347}
{
    \zA{袭人忙连声答应，问：“做什么？”}
}
{
    \eA{[English source not present in the checked-in Bencraft/Joly files.]}
}
{
}

\Verse{348}
{
    \zA{宝玉因要茶吃。}
}
{
    \eA{[English source not present in the checked-in Bencraft/Joly files.]}
}
{
}

\Verse{349}
{
    \zA{袭 人倒了茶来，宝玉乃叹道：“我近来叫惯了他，却忘了是你。”}
}
{
    \eA{[English source not present in the checked-in Bencraft/Joly files.]}
}
{
}

\Verse{350}
{
    \zA{袭人笑道：“他乍 来，你也曾睡梦中叫我，以后才改了的。”}
}
{
    \eA{[English source not present in the checked-in Bencraft/Joly files.]}
}
{
}

\Verse{351}
{
    \zA{说着，大家又睡下。}
}
{
    \eA{[English source not present in the checked-in Bencraft/Joly files.]}
}
{
}

\Verse{352}
{
    \zA{宝玉又翻转了一个 更次。}
}
{
    \eA{[English source not present in the checked-in Bencraft/Joly files.]}
}
{
}

\Verse{353}
{
    \zA{至五更方睡去时，只见晴雯从外走来，仍是往日行景，进来向宝玉道：“你 们好生过罢。}
}
{
    \eA{[English source not present in the checked-in Bencraft/Joly files.]}
}
{
}

\Verse{354}
{
    \zA{我从此就别过了！”}
}
{
    \eA{[English source not present in the checked-in Bencraft/Joly files.]}
}
{
}

\Verse{355}
{
    \zA{说毕，翻身就走。}
}
{
    \eA{[English source not present in the checked-in Bencraft/Joly files.]}
}
{
}

\Verse{356}
{
    \zA{宝玉忙叫时，又将袭人叫醒。}
}
{
    \eA{[English source not present in the checked-in Bencraft/Joly files.]}
}
{
}

\Verse{357}
{
    \zA{袭人还只当他惯了口乱叫，却见宝玉哭了，说道：“晴雯死了！”}
}
{
    \eA{[English source not present in the checked-in Bencraft/Joly files.]}
}
{
}

\Verse{358}
{
    \zA{袭人笑道：“这 是那里的话？}
}
{
    \eA{[English source not present in the checked-in Bencraft/Joly files.]}
}
{
}

\Verse{359}
{
    \zA{叫人听着什么意思。”}
}
{
    \eA{[English source not present in the checked-in Bencraft/Joly files.]}
}
{
}

\Verse{360}
{
    \zA{宝玉那里肯听？}
}
{
    \eA{[English source not present in the checked-in Bencraft/Joly files.]}
}
{
}

\Verse{361}
{
    \zA{恨不得一时亮了就遣人去问信。}
}
{
    \eA{[English source not present in the checked-in Bencraft/Joly files.]}
}
{
}

\Verse{362}
{
    \zA{及至亮时，就有王夫人房里小丫头叫开前角门，传王夫人的话：“‘即时叫起 宝玉，快洗脸换了衣裳来。}
}
{
    \eA{[English source not present in the checked-in Bencraft/Joly files.]}
}
{
}

\Verse{363}
{
    \zA{因今儿有人请老爷赏秋菊，老爷因喜欢他前儿做的诗好， 故此要带了他们去。’}
}
{
    \eA{[English source not present in the checked-in Bencraft/Joly files.]}
}
{
}

\Verse{364}
{
    \zA{这都是太太的话，你们快告诉去，立逼他快来，老爷在上屋 里等他们吃面茶呢。}
}
{
    \eA{[English source not present in the checked-in Bencraft/Joly files.]}
}
{
}

\Verse{365}
{
    \zA{环哥儿早来了。}
}
{
    \eA{[English source not present in the checked-in Bencraft/Joly files.]}
}
{
}

\Verse{366}
{
    \zA{快快儿的去罢。}
}
{
    \eA{[English source not present in the checked-in Bencraft/Joly files.]}
}
{
}

\Verse{367}
{
    \zA{我去叫兰哥儿去了。”}
}
{
    \eA{[English source not present in the checked-in Bencraft/Joly files.]}
}
{
}

\Verse{368}
{
    \zA{里面的 婆子听一句，应一句，一面扣着钮子，一面开门。}
}
{
    \eA{[English source not present in the checked-in Bencraft/Joly files.]}
}
{
}

\Verse{369}
{
    \zA{袭人听得叩门，便知有事，一面 命人问时，自己已起来了。}
}
{
    \eA{[English source not present in the checked-in Bencraft/Joly files.]}
}
{
}

\Verse{370}
{
    \zA{听得这话，忙催人来舀了洗脸水，催宝玉起来梳洗，他 自去取衣。}
}
{
    \eA{[English source not present in the checked-in Bencraft/Joly files.]}
}
{
}

\Verse{371}
{
    \zA{因思跟贾政出门，便不肯拿出十分出色的新鲜衣服来，只拣那三等成色 的来。}
}
{
    \eA{[English source not present in the checked-in Bencraft/Joly files.]}
}
{
}

\Verse{372}
{
    \zA{宝玉此时已无法，只得忙忙前来。}
}
{
    \eA{[English source not present in the checked-in Bencraft/Joly files.]}
}
{
}

\Verse{373}
{
    \zA{果然贾政在那里吃茶，十分喜悦。}
}
{
    \eA{[English source not present in the checked-in Bencraft/Joly files.]}
}
{
}

\Verse{374}
{
    \zA{宝玉请 了早安，贾环贾兰二人也都见过，贾政命坐吃茶，向环兰二人道：“宝玉读书，不 及你两个；论题联、和诗这种聪明，你们皆不及他。}
}
{
    \eA{[English source not present in the checked-in Bencraft/Joly files.]}
}
{
}

\Verse{375}
{
    \zA{今日此去，未免叫你们做诗， 宝玉须随便助他们两个。”}
}
{
    \eA{[English source not present in the checked-in Bencraft/Joly files.]}
}
{
}

\Verse{376}
{
    \zA{王夫人自来不曾听见这等考语，真是意外之喜。}
}
{
    \eA{[English source not present in the checked-in Bencraft/Joly files.]}
}
{
}

\Verse{377}
{
    \zA{一时候他父子去了，方欲过贾 母那边来时，就有芳官等三个干娘走来，回说：“芳官自前日蒙太太的恩典赏出来
        了，他就疯了似的，茶饭都不吃，勾引上藕官蕊官，三个人寻死觅活，只要铰了头 发做尼姑去。}
}
{
    \eA{[English source not present in the checked-in Bencraft/Joly files.]}
}
{
}

\Verse{378}
{
    \zA{我只当是小孩子家，一时出去不惯，也是有的，不过隔两日就好了， 谁知越闹越凶，打骂着也不怕。}
}
{
    \eA{[English source not present in the checked-in Bencraft/Joly files.]}
}
{
}

\Verse{379}
{
    \zA{实在没法，所以来求太太，或是依他们去做尼姑去， 或教导他们一顿，赏给别人做女孩儿去罢。}
}
{
    \eA{[English source not present in the checked-in Bencraft/Joly files.]}
}
{
}

\Verse{380}
{
    \zA{我们没这福。”}
}
{
    \eA{[English source not present in the checked-in Bencraft/Joly files.]}
}
{
}

\Verse{381}
{
    \zA{王夫人听了，道：“胡 说!那里由得他们起来？}
}
{
    \eA{[English source not present in the checked-in Bencraft/Joly files.]}
}
{
}

\Verse{382}
{
    \zA{佛门也是轻易进去的么？}
}
{
    \eA{[English source not present in the checked-in Bencraft/Joly files.]}
}
{
}

\Verse{383}
{
    \zA{每人打一顿给他们，看还闹不闹！”}
}
{
    \eA{[English source not present in the checked-in Bencraft/Joly files.]}
}
{
}

\Verse{384}
{
    \zA{当下因八月十五日各庙内上供去，皆有各庙内的尼姑来送供尖，因曾留下水月庵的 智通与地藏庵的圆信住下未回，听得此信，就想拐两个女孩子去做活使唤。}
}
{
    \eA{[English source not present in the checked-in Bencraft/Joly files.]}
}
{
}

\Verse{385}
{
    \zA{都向王 夫人说：“府上到底是善人家。}
}
{
    \eA{[English source not present in the checked-in Bencraft/Joly files.]}
}
{
}

\Verse{386}
{
    \zA{因太太好善，所以感应得这些小姑娘们皆如此。}
}
{
    \eA{[English source not present in the checked-in Bencraft/Joly files.]}
}
{
}

\Verse{387}
{
    \zA{虽 然说‘佛门容易难上’，也要知道‘佛法平等’，我佛立愿，原度一切众生。}
}
{
    \eA{[English source not present in the checked-in Bencraft/Joly files.]}
}
{
}

\Verse{388}
{
    \zA{如今 两三个姑娘既然无父母，家乡又远，他们既经了这富贵，又想从小命苦，入了风流 行次，将来知道终身怎么样？}
}
{
    \eA{[English source not present in the checked-in Bencraft/Joly files.]}
}
{
}

\Verse{389}
{
    \zA{所以‘苦海回头’，立意出家，修修来世，也是他们 的高意。}
}
{
    \eA{[English source not present in the checked-in Bencraft/Joly files.]}
}
{
}

\Verse{390}
{
    \zA{太太倒不要阻了善念。”}
}
{
    \eA{[English source not present in the checked-in Bencraft/Joly files.]}
}
{
}

\Verse{391}
{
    \zA{王夫人原是个善人，起先听见这话，谅系小孩子 不遂心的话，将来熬不得清净，反致获罪。}
}
{
    \eA{[English source not present in the checked-in Bencraft/Joly files.]}
}
{
}

\Verse{392}
{
    \zA{今听了这两个拐子的话，大近情理。}
}
{
    \eA{[English source not present in the checked-in Bencraft/Joly files.]}
}
{
}

\Verse{393}
{
    \zA{且 近日家中多故，又有邢夫人遣人过来知会，明日接迎春家去住两日，以备人家相看； 且又有官媒来求说探春等，心绪正烦，那里着意在这些小事？}
}
{
    \eA{[English source not present in the checked-in Bencraft/Joly files.]}
}
{
}

\Verse{394}
{
    \zA{既听此言，便笑答道： “你两个既这等说，你们就带了做徒弟去，如何？”}
}
{
    \eA{[English source not present in the checked-in Bencraft/Joly files.]}
}
{
}

\Verse{395}
{
    \zA{二姑子听了，念一声佛，道： “善哉，善哉!若如此，可是老人家的阴功不小。”}
}
{
    \eA{[English source not present in the checked-in Bencraft/Joly files.]}
}
{
}

\Verse{396}
{
    \zA{说毕便稽首拜谢。}
}
{
    \eA{[English source not present in the checked-in Bencraft/Joly files.]}
}
{
}

\Verse{397}
{
    \zA{王夫人道： “既这样，你们问他去。}
}
{
    \eA{[English source not present in the checked-in Bencraft/Joly files.]}
}
{
}

\Verse{398}
{
    \zA{若果真心，即上来当着我拜了师父去罢。”}
}
{
    \eA{[English source not present in the checked-in Bencraft/Joly files.]}
}
{
}

\Verse{399}
{
    \zA{这三个女人听了出去，果然将他三人带来。}
}
{
    \eA{[English source not present in the checked-in Bencraft/Joly files.]}
}
{
}

\Verse{400}
{
    \zA{王夫人问之再三，他三人已立定主 意，遂与两个姑子叩了头，又拜辞了王夫人。}
}
{
    \eA{[English source not present in the checked-in Bencraft/Joly files.]}
}
{
}

\Verse{401}
{
    \zA{王夫人见他们意皆决断，知不可强了， 反倒伤心可怜，忙命人来取了些东西来赏了他们，又送了两个姑子些礼物。}
}
{
    \eA{[English source not present in the checked-in Bencraft/Joly files.]}
}
{
}

\Verse{402}
{
    \zA{从此芳 官跟了水月庵的智通，蕊官藕官二人跟了地藏庵圆信，各自出家去了。}
}
{
    \eA{[English source not present in the checked-in Bencraft/Joly files.]}
}
{
}

\Verse{403}
{
    \zA{要知后事，下回分解。}
}
{
    \eA{[English source not present in the checked-in Bencraft/Joly files.]}
}
{
}

\Verse{404}
{
    \zA{(本章完)}
}
{
    \eA{[English source not present in the checked-in Bencraft/Joly files.]}
}
{
}
